\documentclass[11pt,a4paper]{article}
\usepackage[utf8]{inputenc}
\usepackage[spanish]{babel}
\usepackage{amsmath}
\usepackage{amsfonts}
\usepackage{amssymb}
\usepackage{geometry}
\usepackage{enumerate}
\geometry{top=2.5cm, bottom=2.5cm, left=2.5cm, right=2.5cm}

\begin{document}

\begin{center}
    \textbf{\Large SEGUNDO PARCIAL DE FÍSICA BÁSICA} \\[0.2cm]
    \textbf{\large VERSIÓN A} \\[0.4cm]
\end{center}

\noindent \textbf{Estudiante:} \hrulefill \ \textbf{Matrícula:} \underline{\hspace{3cm}} \\[0.2cm]
\textbf{Fecha:} \underline{\hspace{4cm}} \ \textbf{Calificación:} \underline{\hspace{1.5cm}} / 8.0 puntos \\

\hrule
\section*{I. Teoría (Valor: 4.0 puntos - 0.5 puntos cada una)}
\textit{Selecciona la opción correcta en cada una de las siguientes preguntas. No debe haber tachaduras ni enmiendas.}

\begin{enumerate}
    \item \textbf{A diferencia de la masa, el peso de un objeto:}
    \begin{enumerate}[a)]
        \item Es una propiedad intrínseca que no varía con la gravedad local.
        \item Es la fuerza de atracción gravitatoria que ejerce la Tierra sobre él ($w = m \cdot g$).
        \item Se mide en kilogramos en el Sistema Internacional.
        \item Es una cantidad escalar pura sin dirección ni sentido.
    \end{enumerate}

    \item \textbf{De acuerdo con la Tercera Ley de Newton (Acción-Reacción), si un cuerpo A ejerce una fuerza $\vec{F}_{AB}$ sobre un cuerpo B, entonces:}
    \begin{enumerate}[a)]
        \item El cuerpo B ejerce sobre A una fuerza de igual magnitud, misma dirección y sentido contrario ($\vec{F}_{BA} = -\vec{F}_{AB}$).
        \item El cuerpo B ejerce sobre A una fuerza de mayor magnitud en el mismo sentido.
        \item La fuerza del cuerpo B sobre A es nula porque A inició el contacto.
        \item El cuerpo B ejerce sobre A una fuerza perpendicular a la original.
    \end{enumerate}

    \item \textbf{En un gráfico de posición en función del tiempo ($x$ vs $t$), la velocidad instantánea en un tiempo determinado equivale a:}
    \begin{enumerate}[a)]
        \item El área total bajo la curva en ese intervalo.
        \item La pendiente de la recta secante que corta la curva.
        \item La pendiente de la recta tangente a la curva en ese punto.
        \item La longitud del vector de aceleración.
    \end{enumerate}

    \item \textbf{Despreciando la resistencia del aire, el ángulo de disparo respecto a la horizontal que produce el máximo alcance horizontal para un proyectil es:}
    \begin{enumerate}[a)]
        \item $30^\circ$
        \item $45^\circ$
        \item $60^\circ$
        \item $90^\circ$
    \end{enumerate}

    \item \textbf{La fuerza es una magnitud física de carácter:}
    \begin{enumerate}[a)]
        \item Escalar.
        \item Vectorial.
        \item Adimensional.
        \item Tensorial de segundo orden.
    \end{enumerate}

    \item \textbf{La propiedad física de un cuerpo que mide cuantitativamente su inercia (resistencia al cambio en su estado de movimiento) es la:}
    \begin{enumerate}[a)]
        \item Aceleración.
        \item Fuerza normal.
        \item Masa.
        \item Fricción.
    \end{enumerate}

    \item \textbf{De acuerdo con la Primera Ley de Newton, si la fuerza neta que actúa sobre un cuerpo es nula, el cuerpo necesariamente:}
    \begin{enumerate}[a)]
        \item Debe estar en reposo absoluto.
        \item Permanecerá en reposo o en movimiento rectilíneo uniforme.
        \item Experimentará un aumento constante de su velocidad.
        \item Cambiará continuamente su dirección de movimiento.
    \end{enumerate}

    \item \textbf{La Segunda Ley de Newton establece matemáticamente que la aceleración de un cuerpo es:}
    \begin{enumerate}[a)]
        \item Directamente proporcional a su masa e inversamente proporcional a la fuerza neta.
        \item Directamente proporcional a la fuerza neta e inversamente proporcional a su masa.
        \item Independiente de la masa y proporcional al peso.
        \item Proporcional al cuadrado de la fuerza neta aplicada.
    \end{enumerate}

\end{enumerate}

\hrule
\section*{II. Problemas Prácticos (Valor: 4.0 puntos - 1.0 punto cada uno)}
\textit{Resuelve cada uno de los problemas y selecciona la respuesta correcta. Muestra tus cálculos breves.}

\begin{enumerate}
    \item \textbf{Un automóvil que viaja en línea recta a 90.0 \text{ km/h} (25.0 \text{ m/s}) pisa los frenos durante 4.00 \text{ s} reduciendo su velocidad de forma uniforme hasta 54.0 \text{ km/h} (15.0 \text{ m/s}). Calcule la distancia recorrida ($\Delta x$) durante el frenado:}
    \begin{enumerate}[a)]
        \item $71.50 \text{ m}$
        \item $92.00 \text{ m}$
        \item $80.00 \text{ m}$
        \item $100.00 \text{ m}$
    \end{enumerate}

    \item \textbf{Una caja se desliza por una rampa inclinada con una aceleración constante de 2.44 \text{ m/s}^2$ partiendo con una velocidad inicial de 1.13 \text{ m/s}. Calcule el desplazamiento de la caja al cabo de 2.00 \text{ s}$:}
    \begin{enumerate}[a)]
        \item $2.26 \text{ m}$
        \item $12.14 \text{ m}$
        \item $3.64 \text{ m}$
        \item $7.14 \text{ m}$
    \end{enumerate}

    \item \textbf{Calcule la magnitud de la fuerza neta que actúa sobre un bloque de masa $m = 3.27 \text{ kg}$ si este se mueve con una aceleración constante de 0.95 \text{ m/s}^2$:}
    \begin{enumerate}[a)]
        \item $2.31 \text{ N}$
        \item $3.44 \text{ N}$
        \item $3.11 \text{ N}$
        \item $4.61 \text{ N}$
    \end{enumerate}

    \item \textbf{Determine la masa de un cuerpo que se mueve con una aceleración de 2.00 \text{ m/s}^2$ sabiendo que la fuerza neta que actúa sobre él es $\vec{F}_{\text{neta}} = (6.00 \text{ N}, 8.00 \text{ N})$:}
    \begin{enumerate}[a)]
        \item $4.20 \text{ kg}$
        \item $5.00 \text{ kg}$
        \item $6.50 \text{ kg}$
        \item $20.00 \text{ kg}$
    \end{enumerate}

\end{enumerate}

\end{document}
